Dáta spracované do formátu BAM obsahujú viacero sekvencií DNA, ktoré sa navzájom líšia dĺžkou aj obsahom.
Každý riadok reprezentuje jedno čítanie z nanopórovej sekvenácie.
Dĺžka DNA čítaní sa pohybuje približne od 100 nukleotidov až po 1~000~000 nukleotidov.

Jedným z možných dôvodov výskytu veľmi dlhých čítaní je, že zariadenie môže postupne načítať viac DNA sekvencií za sebou, pričom nástroj Dorado ich neoddelí a spracuje ich ako jednu sekvenciu.
Zároveň možno predpokladať, že v rámci jedného čítania sa tá istá sekvencia môže vyskytovať opakovane.
Môže to súvisieť s biosyntézou DNA, pri ktorej dochádza k replikácii a vzniku viacerých kópií tej istej sekvencie.

Tá istá sekvencia môže byť načítaná viackrát, avšak s odlišnými chybami, keďže nanopórové sekvenovanie je chybovejšie.
Okrem toho sa sekvencia môže čítať v opačnej orientácii, pretože molekula vstupuje do nanopóru z rôznych koncov.
Sekvenciu preto možno získať v smere 3$'$\textrightarrow5$'$ alebo 5$'$\textrightarrow3$'$.

V dôsledku biosyntézy vzniká aj komplementárna sekvencia DNA, v ktorej sú bázy nahradené podľa pravidiel párovania (A~\textrightarrow~T, T~\textrightarrow~A, G~\textrightarrow~C, C~\textrightarrow~G).
Aj komplementárna sekvencia môže byť čítaná v oboch smeroch.
Pre každú pôvodnú sekvenciu sa teda môžu vyskytnúť až štyri varianty: pôvodná sekvencia v smere 3$'$\textrightarrow5$'$ a 5$'$\textrightarrow3$'$, a jej komplement v smere 3$'$\textrightarrow5$'$ a 5$'$\textrightarrow3$'$.

Ďalším javom je fragmentácia DNA, pri ktorej vznikajú kratšie sekvencie zodpovedajúce iba časti pôvodnej sekvencie.
Ich dĺžky sa líšia podľa miesta pretrhnutia.
Zároveň sa môže stať, že niektoré sekvencie sa replikujú častejšie ako iné.


\begin{comment}
Napriek uvedeným komplikáciám som sa pokúsil všetky DNA čítania klastrovať.
Identifikoval som klaster, ktorý obsahoval viac ako 37~000 sekvencií DNA.
Spočiatku sa zdalo, že ide o opakované čítanie jednej sekvencie, čo by mohlo viesť ku kvalitným dátam.
Podrobnejšia analýza však ukázala, že sekvencie v tomto klastri mali spoločný iba úvodný úsek s dĺžkou 43 nukleotidov.

Následne som sa pokúsil klastrovať úseky dĺžky 37~000 nukleotidov za daným primerom, no bez úspechu.
Preto som primer vyhľadával v každom čítaní, nielen na začiatku sekvencie, ale aj v jej ďalších častiach.
Následne som opätovne použil nástroj Clover, tentokrát iba na zvyšné časti sekvencií po odstránení DNA primerov.
\end{comment}

Napriek uvedeným komplikáciám som predpokladal, že v dlhých čítaniach bude možné identifikovať dostatočne dlhé opakujúce sa úseky, ktoré reprezentujú tú istú pôvodnú DNA sekvenciu.
Na základe tohto predpokladu som identifikoval klaster obsahujúci viac ako 37~000 sekvencií DNA.
Spočiatku sa zdalo, že ide o opakované čítanie jednej sekvencie, čo by mohlo viesť k získaniu kvalitného konsenzu.
Podrobnejšia analýza však ukázala, že spoločná zhoda medzi sekvenciami bola obmedzená iba na úvodný úsek s dĺžkou 43 nukleotidov.

Keďže 43 nukleotidov predstavuje príliš krátky úsek na spoľahlivú reprezentáciu celého DNA reťazca no zaroven je priliš dlhý na náhodne sa vyskytujúci úsek, vznikla
, pôvodná hypotéza o prítomnosti dostatočne dlhého spoločného jadra nebola potvrdená.
Na základe tohto zistenia bola formulovaná nová hypotéza: spoločný úvodný úsek pravdepodobne zodpovedá DNA primeru prítomnému na začiatku sekvenovaných fragmentov.

Nasledujúcim krokom preto bolo lokálne hľadanie primeru v jednotlivých čítaniach.
Primer sa nevyhľadával iba na začiatku reťazca, ale v celom rozsahu každého čítania, aby sa zachytili aj prípady jeho posunu alebo výskytu vo vnútri sekvencie.

%AK najdeny primer => dna za prijmerom je spravne citanie a nie je to komplementarny retazec alebo retazec citany v opacnom smere

%\begin{lstlisting}[language=Python, caption={Clustering Data}, label={lst:python-clustering}] 
    tuples = m.data
    data_sorted = sorted(tuples, key=lambda x: x[1])
    groups = []

    for k, g in itertools.groupby(data_sorted, key=lambda x: x[1]):
        groups.append([k] + [a[0] for a in g])
    return groups
\end{lstlisting}